\chapter{Creio na ressurreição da carne}
\chaptermark{Creio na ressurreição da carne}

Este capítulo conclui o terceiro artigo do Símbolo Apostólico, professando a fé na ressurreição da carne e na vida eterna. Desdobramos os parágrafos correspondentes do Catecismo da Igreja Católica (CIC 988--1019), culminando o Credo na vitória sobre a morte por Cristo, com esperança na ressurreição corporal e vida com Deus. Para adultos em processo de iniciação cristã pelo RICA, essa profissão ilumina o Batismo pascal como germe de imortalidade.

\section{A Ressurreição dos Mortos (\S 988--1004)}

O Credo cristão culmina na proclamação da ressurreição dos mortos no último dia e na vida eterna, professando fé na ação criadora, salvífica e santificadora de Deus.

Cremos firmemente que, como Cristo ressuscitou verdadeiramente, os justos viverão para sempre com Ele, raised up pelo Pai, Filho e Espírito Santo no último dia.

``Ressurreição da carne'' significa que o corpo mortal reviverá, além da imortalidade da alma; ``carne'' indica o homem em fraqueza mortal.3

Essa crença é essencial desde os inícios cristãos: sem ela, a fé é vã; Cristo é primícias dos adormecidos.

Revelação progressiva: esperança no Criador e Aliança com Abraão; mártires macabeus confessam ressurreição pela fidelidade às leis divinas.

Jesus ensina firmemente a ressurreição aos saduceus, ligando-a à sua pessoa: ``Eu sou a Ressurreição e a Vida'', restaurando mortos como sinal.

Para adultos a partir dos 20 anos no RICA, meditar a ressurreição no rito de iniciação pascal desperta esperança escatológica, unindo catecúmenos à vitória de Cristo na Vigília.

\section{Morrer e Ressuscitar com Cristo; A Morte (\S 1005--1019)}

Para ressuscitar com Cristo, devemos morrer com Ele: na morte, alma separa-se do corpo, aguardando reunião gloriosa.

A morte é natural mas salário do pecado; para os que morrem em graça, participa da morte do Senhor rumo à Ressurreição.

Fim da peregrinação terrena, urge santificar o tempo limitado; não há reencarnação, mas juízo único após uma só vida.

Cristo transforma a morte: obediência filial converte maldição em bênção; cristãos veem-na como ganho, completando união batismal.

Liturgia exprime: ``vida se muda, não se acaba''; prepara-se com litania dos santos, Maria e São José, patrono da boa morte.

Na catequese para adultos não iniciados, contemplar morte cristã no tempo quaresmal fortalece neófitos na Penitência, rumo à Eucaristia como penhor de ressurreição.

\section{Conclusão}

Professar ressurreição da carne e vida eterna afirma triunfo trinitário sobre morte. Adultos no RICA, vivam Batismo como já ressuscitados, ansiando Parusia.